\documentclass[11pt]{article}
\usepackage[a4paper,margin=27mm]{geometry}
\usepackage[T1]{fontenc}
\usepackage{lmodern,microtype,amsmath,amssymb,amsthm,mathtools,enumitem}
\usepackage[colorlinks=true,linkcolor=blue,citecolor=blue]{hyperref}
\newtheorem{theorem}{Theorem}
\newtheorem{lemma}[theorem]{Lemma}
\newtheorem{corollary}[theorem]{Corollary}
\newtheorem{proposition}[theorem]{Proposition}
\theoremstyle{remark}
\newtheorem{remark}[theorem]{Remark}
\newcommand{\R}{\mathbb R}
\newcommand{\one}{\mathbf 1}
\newcommand{\supp}{\operatorname{supp}}
\newcommand{\sd}{\mathbin{\triangle}}
\newcommand{\dH}{d_{\mathrm H}}
\newcommand{\lo}{\underline P}
\newcommand{\hi}{\overline P}
\title{Full Coverage by Homothetic Credal Families}
\author{}
\date{8 October 2026}
\begin{document}
\maketitle

\noindent
Every finite-dimensional credal set belongs to a continuous family satisfying
face support and compatibility of covers. This includes the vacuous set and
all other vertex-containing credal sets. Combined with the ordinal refinement
of \cite{OR}, the construction gives rankings on all measurable models which
strictly respect background error inclusion and collectively satisfy
\textbf{(E)}, \textbf{(LA)}, \textbf{(A)}, and \textbf{(S)}. With coherence
understood modulo null sets, they also satisfy \textbf{(C)}.

The observation is that common covers order both endpoints of a marginal
probability interval. Indexing by a fixed weighted average of those endpoints
therefore suffices. Upper probabilities alone need not order the indices.

\section{Common covers order both endpoints}

Fix $n\ge2$, $\Delta=\{p\in\R^n:p\ge0,\ \sum_i p_i=1\}$, and any
finite probability measure $\mu$ equivalent to Lebesgue measure. All inclusions
and equalities between models below are understood modulo null sets where
indicated. Models are all Lebesgue-measurable subsets of $\R^n$.

For a nonempty closed convex $K\subseteq\Delta$, put
\[
 \varphi_K(g)=\min_{m\in K}m\cdot g,\qquad
 D_K=\{g:\varphi_K(g)>0\}\cup\{g:g\ge0,\ g\ne0\},
\]
and define the marginal endpoints
\[
 \lo_K(i)=\min_{m\in K}m_i,\qquad
 \hi_K(i)=\max_{m\in K}m_i.
\]
Let $G_i=\{g:g_i>0\}$, $E_i(X)=X\sd G_i$, and
\[
 X\le_iY\quad\Longleftrightarrow\quad
 E_i(X)\subseteq E_i(Y)\text{ a.e.}
\]
Its strict part is $<_i$. Let $W_{ij}=\{g:g_i>0\ge g_j\}$ for $i\ne j$.

\begin{lemma}\label{lem:endpoints}
Suppose $i\ne j$, $X\le_iD_K$, and $X\le_jD_H$. Then
\[
 \lo_K(i)\le\lo_H(i),\qquad
 \hi_K(i)\le\hi_H(i).
 \tag{1}\label{eq:endpoints}
\]
\end{lemma}
\begin{proof}
On $W_{ij}$, the two background comparisons give
\[
 D_K\cap W_{ij}\subseteq X\cap W_{ij}\subseteq D_H\cap W_{ij}
 \quad\text{a.e.}
\]
Exchanging $i$ and $j$ gives $D_H\cap W_{ji}\subseteq D_K$ a.e.

If $\lo_K(i)>\lo_H(i)$, choose $t$ strictly between them. Then $0<t<1$,
and $g=e_i-t\one$ satisfies
\[
 g_i>0>g_j,\qquad
 \varphi_K(g)=\lo_K(i)-t>0>\lo_H(i)-t=\varphi_H(g).
\]
The strict inequalities persist on an open neighbourhood, contradicting
$D_K\cap W_{ij}\subseteq D_H$ a.e.

If $\hi_K(i)>\hi_H(i)$, choose $t$ strictly between them. Now
$g=t\one-e_i$ satisfies
\[
 g_j>0>g_i,\qquad
 \varphi_H(g)=t-\hi_H(i)>0>t-\hi_K(i)=\varphi_K(g).
\]
Again an open neighbourhood contradicts $D_H\cap W_{ji}\subseteq D_K$ a.e.
\end{proof}

\begin{corollary}\label{cor:cc}
Suppose a family $p\mapsto K_p$ satisfies, for some fixed $0<\lambda<1$,
\[
 (1-\lambda)\lo_{K_p}(i)+\lambda\hi_{K_p}(i)=p_i
 \quad\text{for every }p,i.
 \tag{2}\label{eq:index}
\]
Then it satisfies compatibility of covers:
\[
 X\le_iD_{K_p},\quad X\le_jD_{K_q},\quad i\ne j
 \quad\Longrightarrow\quad p_i\le q_i.
 \tag{CC}\label{eq:cc}
\]
\end{corollary}
\begin{proof}
Take the weighted average in \eqref{eq:endpoints}.
\end{proof}

\section{A family through an arbitrary target}

Let $M\subseteq\Delta$ be any nonempty closed convex set. Put
\[
 \ell_i=\lo_M(i),\qquad u_i=\hi_M(i),\qquad
 L=\sum_i\ell_i,\qquad U=\sum_i u_i.
\]
If $M$ is not a singleton, then $L<1<U$: equality in either bound would
force every member of $M$ to equal the same endpoint vector. Define
\[
 \lambda=\frac{1-L}{U-L},\qquad
 r=(1-\lambda)\ell+\lambda u,
 \qquad d=r-\ell=\lambda(u-\ell).
 \tag{3}\label{eq:anchor}
\]
Then $0<\lambda<1$ and $r\in\Delta$. If $M=\{m\}$, instead take
$\lambda=1/2$, $r=m$, and $d=0$. The anchor $r$ need not belong to $M$;
only $r\in\Delta$ is needed.

For $p\in\Delta$, set
\[
 \theta(p)=\min\left\{1,\ \min_{i:d_i>0}\frac{p_i}{d_i}\right\},
 \qquad
 \boxed{\quad K_p=p+\theta(p)(M-r).\quad}
 \tag{4}\label{eq:family}
\]
The minimum over an empty set is $+\infty$, so for singleton targets
$\theta\equiv1$ and $K_p=\{p\}$.

\begin{theorem}[Full family coverage]\label{thm:family}
For every target $M$, the family \eqref{eq:family} has the following properties.
\begin{enumerate}[label=(\alph*)]
\item Each $K_p$ is a nonempty closed convex subset of $\Delta$, and $K_r=M$.
\item The endpoint identity \eqref{eq:index} holds. Hence (CC) holds.
\item $K_p\subseteq\Delta_{\supp p}$ for every $p$.
\item $p\mapsto K_p$ is Hausdorff-Lipschitz. In particular
$\mu(D_{K_p}\sd D_{K_q})\to0$ as $q\to p$.
\end{enumerate}
No interiority or vertex exclusion assumption on $M$ is required.
\end{theorem}
\begin{proof}
An affine image of the compact convex set $M$ is nonempty, compact and convex.
Every point of $K_p$ has coordinate sum one, and its minimum $i$th coordinate is
\[
 \lo_{K_p}(i)=p_i+\theta(p)(\ell_i-r_i)
             =p_i-\theta(p)d_i\ge0.
 \tag{5}\label{eq:lower}
\]
Thus $K_p\subseteq\Delta$. Since $r_i=d_i+\ell_i\ge d_i$, we have
$\theta(r)=1$, giving $K_r=M$. This proves (a).

The maximum coordinate is
\[
 \hi_{K_p}(i)=p_i+\theta(p)(u_i-r_i)
             =p_i+\theta(p)(1-\lambda)(u_i-\ell_i).
 \tag{6}\label{eq:upper}
\]
Combining \eqref{eq:lower}, \eqref{eq:upper}, and
$d_i=\lambda(u_i-\ell_i)$ gives \eqref{eq:index}; Corollary~\ref{cor:cc}
proves (CC).

For (c), let $p_i=0$. If $d_i>0$, then $\theta(p)=0$ and $K_p=\{p\}$.
If $d_i=0$, then $u_i=\ell_i=r_i$, so every $m\in M$ has $m_i=r_i$ and
every member of $K_p$ has coordinate $p_i+\theta(p)(m_i-r_i)=0$.

For (d), $\theta$ is the minimum of finitely many affine functions and the
constant one, hence Lipschitz. Pairing the images of the same $m\in M$ gives
\[
 \dH(K_p,K_q)
 \le\|p-q\|_1+|\theta(p)-\theta(q)|\max_{m\in M}\|m-r\|_1
 \le\|p-q\|_1+2|\theta(p)-\theta(q)|.
\]
Hausdorff convergence makes the lower expectations converge pointwise. Their
zero sets are Lebesgue-null, since constant additivity gives exactly one zero
on every line parallel to $\one$. Desirable sets differ from their strict
positive regions only on null sets. Dominated convergence therefore gives the
stated continuity in symmetric-difference measure.
\end{proof}

\begin{remark}
The family also satisfies $\hi_{K_p}(i)\le p_i/\lambda$. For $d_i>0$ this
follows from $\theta(p)d_i\le p_i$ and \eqref{eq:upper}; for $d_i=0$ the upper
probability is exactly $p_i$. The ordinal refinement does not need this bound.
\end{remark}

\section{The ordinal coverage theorem}

For clarity, recall the construction from \cite{OR}. For a fixed family let
\[
 \mathcal P_i(X)=\{p:X\le_iD_{K_p}\},\qquad
 T(X)=\{i:\mathcal P_i(X)\ne\varnothing\},
\]
\[
 a_i(X)=\max_{p\in\mathcal P_i(X)}p_i\quad(i\in T(X)),\qquad
 s(X)=1-\sum_{i\in T(X)}a_i(X).
\]
The maxima exist by continuity. The ordinal refinement uses
\[
 b_i(X)=\inf_{Y\le_iX}\left(a_i(Y)+\frac{s(Y)}{2n}\right)
 \quad(i\in T(X)).
 \tag{7}\label{eq:envelope}
\]
For covered worlds, compare $(a_i(X),b_i(X))$ coordinatewise, with larger
coordinates better; assign $(-1,-1)$ at uncovered worlds. Inside an equal-pair
class use background inclusion, except that all models at a nonnegative
degenerate pair $(t,t)$ tie. This defines preorders, generally with
incomparabilities.

Under (CC), face support and continuity, \cite[Theorem~4]{OR} proves that these
preorders strictly respect background inclusion, that every selected model
$D_{K_p}$ is undominated and competitive on $\supp p$, and that
\[
 \{p:D_{K_p}\preceq_iX\}=
 \begin{cases}
 \{p:p_i\ge b_i(X)\}&i\in T(X),\\
 \Delta&i\notin T(X).
 \end{cases}
 \tag{8}\label{eq:sections}
\]
The undominated models are exactly the null modifications of family members;
every other model is strictly dominated at every world by a family member.
The essential compactness argument proving positivity of the envelope is given
in full in \cite{OR}; no quantitative estimates on error measures are used.

\begin{theorem}[Full ordinal coverage]\label{thm:coverage}
Fix any finite number $n\ge2$ of worlds. There is a class $\mathfrak L$ of
rankings, each defined on all measurable models, such that:
\begin{enumerate}[label=(\alph*)]
\item every ranking respects background inclusion and makes every
positive-measure background improvement strict;
\item $\mathfrak L$ satisfies \textbf{(E)}, \textbf{(LA)}, \textbf{(A)}, and
\textbf{(S)} of the generalized dominance framework;
\item with coherence understood modulo null sets, every ranking satisfies
\textbf{(C)}; its undominated models are exactly its selected family modulo
null sets, and every other model is strictly dominated at every world by a
selected coherent model.
\end{enumerate}
In particular, every coherent desirable-gamble set, including the vacuous one,
is undominated under at least one ranking in $\mathfrak L$.
\end{theorem}
\begin{proof}
For each credal set $M$, form the family of Theorem~\ref{thm:family} and its
ordinal refinement \eqref{eq:envelope}. Theorem~\ref{thm:family} supplies all
three hypotheses of \cite[Theorem~4]{OR}. That theorem gives (a), (LA), (S),
and (c). Every ranking evaluates every measurable model, so (E) holds.

For (A), every coherent desirable-gamble set $D$ agrees modulo null sets with
$D_M$ for a nonempty closed convex $M\subseteq\Delta$. Indeed, normalize the
nonzero vectors in the dual cone of $\overline D$ to obtain $M$. The dual
vectors are nonnegative because $D$ contains the positive orthant. Cone
duality gives
$\operatorname{int}\overline D=\{g:\min_{m\in M}m\cdot g>0\}$, and a
full-dimensional convex set differs from its interior only on a null set.
The family for this $M$ has $K_r=M$, so $D\approx D_M=D_{K_r}$ is undominated
under its corresponding ranking. The same holds for all models coherent
modulo null sets. Thus every coherent model belongs to some $\mathcal C_L$.
\end{proof}

\begin{remark}[Uniform target index]
If desired, the selection can be reparametrized so that the target is selected
at $u=\one/n$. Let
\[
 \gamma(p)=n\min_i p_i,\qquad
 f(p)=(1-\gamma(p))p+\gamma(p)r.
\]
Then $f$ is continuous, $f(u)=r$, and $\supp f(p)\subseteq\supp p$: on the
boundary $f(p)=p$, while an interior $p$ has full support. Under the same
ranking, choose $u_p=D_{K_{f(p)}}$. It is competitive on $\supp p$, and all
comparison sections remain closed as continuous inverse images of the sets
in \eqref{eq:sections}. Thus $u_u=D_M$, including for targets supported on a
proper face. This reparametrization need not preserve (CC) in the new indices;
it is applied after constructing the ranking and preserves (S) directly.
\end{remark}

\section{The vacuous target and the limit of upper-probability conditions}

For $M=\Delta$ we have $\ell=0$, $u=\one$, $\lambda=1/n$, $r=\one/n$ and
$d=\one/n$. Writing $m(p)=\min_i p_i$, the family reduces to
\[
 \boxed{\quad K_p=p-m(p)\one+n\,m(p)\Delta.\quad}
 \tag{9}\label{eq:vacuous}
\]
Its endpoints are
\[
 \lo_{K_p}(i)=p_i-m(p),\qquad
 \hi_{K_p}(i)=p_i+(n-1)m(p).
\]
It equals $\Delta$ at the uniform index and $\{p\}$ on the boundary. The
endpoint identity is
\[
 p_i=\frac{n-1}{n}\lo_{K_p}(i)+\frac1n\hi_{K_p}(i).
\]
Thus the vacuous target is covered without any special case in the ranking.

The family generally fails the proposed upper-probability condition (U$^*$).
For example, in three worlds take
\[
 p=(.70,.20,.10),\qquad q=(.65,.13,.22).
\]
Their upper-probability vectors are respectively
\[
 (.90,.40,.30),\qquad (.91,.39,.48).
\]
The first upper coordinate increases and the second decreases from $p$ to $q$,
but $p_1>q_1$, violating (U$^*$). There is no common-cover obstruction:
the lower first coordinates are $.60$ and $.52$, so the first required gluing
inclusion already fails by Lemma~\ref{lem:endpoints}.

\begin{proposition}\label{prop:ustar}
Suppose a continuous family of credal sets satisfies face support and
\[
 \hi_{K_p}(i)\le\hi_{K_q}(i),\quad
 \hi_{K_q}(j)\le\hi_{K_p}(j),\quad i\ne j
 \quad\Longrightarrow\quad p_i\le q_i.
 \tag{U$^*$}\label{eq:ustar}
\]
If $K_p$ contains a simplex vertex, then $|\supp p|\le2$. Consequently such
families cannot contain a vertex-containing credal set whose union of supports
has at least three worlds. In particular, they cannot contain $\Delta$ when
$n\ge3$.
\end{proposition}
\begin{proof}
Suppose $e_i\in K_p$. Face support gives $i\in\supp p$. If the support has
at least three worlds, choose distinct $j,k\in\supp p\setminus\{i\}$ and
vary only coordinates $i,j,k$, holding all others fixed. The index $p$ is an
interior point of this two-dimensional triangular slice.

On the slice the map
$x\mapsto(\hi_{K_x}(i),\hi_{K_x}(j))$ is continuous and injective.
Indeed equality of both outputs, with (U$^*$) applied in both directions and
with both choices of the distinguished coordinate, gives equality of $x_i$
and $x_j$. Equality of $x_k$ follows from normalization; all remaining
coordinates were fixed. Invariance of domain therefore makes the image of the
slice's interior open in $\R^2$. But at $p$ its first coordinate is one,
and every image point has first coordinate at most one. This is impossible.
\end{proof}

The endpoint construction avoids this obstruction because (CC) uses actual
common covers. It does not require upper-probability inequalities alone to
determine the index order.

\section{Full coverage by bounded real-valued losses}

Incomparability is not necessary for full coverage. The same families also
meet the hypotheses of the numerical construction in \cite{LR}, once its
assumption (U) is replaced by (CC). This replacement is valid because the
only use of (U) in that construction's main theorem is to prove its
two-world common-cover lemma, whose conclusion is exactly (CC).

In this section fix $\mu$ to be standard Gaussian measure. Let $A$ be a
Hausdorff-Lipschitz constant for the family and set $B=1/\lambda$.
Theorem~\ref{thm:family} and the remark following it give
\[
 \dH(K_p,K_q)\le A\|p-q\|_1,
 \qquad \hi_{K_p}(i)\le Bp_i.
\]
Put $\kappa=8n^2(A+2B)$ and choose $0<\eta<1/(n\kappa)$.
Using the cover sets $\mathcal P_i(X)$ defined above, let
\[
 L_\eta(X,i)=
 \begin{cases}
 \displaystyle\min_{p\in\mathcal P_i(X)}
 \left[1-p_i-\eta\mu\big(E_i(D_{K_p})\setminus E_i(X)\big)\right],
       &\mathcal P_i(X)\ne\varnothing,\\[3pt]
 2+\mu(E_i(X)),&\mathcal P_i(X)=\varnothing.
 \end{cases}
 \tag{10}\label{eq:finite-loss}
\]
The minimum exists: the cover set is compact and the objective is continuous.

\begin{theorem}[Full numerical coverage]\label{thm:numerical}
For each target $M$, the loss \eqref{eq:finite-loss} takes values in
$[0,1]\cup[2,3]$, evaluates every measurable model, and strictly respects
background improvements of positive measure. Moreover:
\begin{enumerate}[label=(\alph*)]
\item $L_\eta(D_{K_p},i)=1-p_i$;
\item the selected models are undominated and competitive on their index
supports, and all selected-model comparison sections are closed;
\item every model outside the selected family modulo null sets is strictly
dominated at every world by a selected model.
\end{enumerate}
The class obtained by letting $M$ range over all credal sets satisfies
\textbf{(E)}, \textbf{(LA)}, \textbf{(A)}, and \textbf{(S)}, as well as
\textbf{(C)} with coherence understood modulo null sets.
\end{theorem}
\begin{proof}
Here is the budget argument underlying \cite{LR}, with (CC) explicit.
For $T=T(X)\ne\varnothing$, choose any covers $p^i\in\mathcal P_i(X)$ and put
\[
 \rho=1-\sum_{i\in T}p^i_i,
 \qquad h_i=\mu\big(E_i(D_{K_{p^i}})\setminus E_i(X)\big).
\]
By (CC), $p^j_k\ge p^k_k$ for distinct $j,k\in T$. Hence
\[
 \rho\ge0,\qquad p^j(T^c)\le\rho,
 \qquad\|p^i-p^j\|_1\le2\rho.
\]
The Gaussian estimates proved in \cite{LR} give
\[
 \mu(D_{K_p}\sd D_{K_q})\le4nA\|p-q\|_1,
\]
and, with $G_T^+=\{g:g_i>0\ (i\in T)\}$ and
$G_T^-=\{g:g_i\le0\ (i\in T)\}$,
\[
 \mu(G_T^+\setminus D_{K_p}),\quad
 \mu(D_{K_p}\cap G_T^-)
 \ \le\ 8n^2B\,p(T^c).
\]
An error removed from $D_{K_{p^j}}$ by $X$ either witnesses disagreement
with another covering model or lies in one of these two face-error regions.
Consequently
\[
 h_j\le\sum_{i\in T\setminus\{j\}}
       \mu(D_{K_{p^i}}\sd D_{K_{p^j}})
       +16n^2B\,p^j(T^c)
 \le\kappa\rho.
\]
It follows that
\[
 \sum_{i\in T}(p^i_i+\eta h_i)
 \le1-(1-n\kappa\eta)\rho\le1.
 \tag{11}\label{eq:numerical-budget}
\]
Choose the covers attaining \eqref{eq:finite-loss} and write
$c_i(X)=1-L_\eta(X,i)$ for $i\in T$. These numbers are nonnegative, and
their sum is at most one, proving the claimed loss range. Equality in the
budget forces $\rho=0$, all chosen covers equal, and all $h_i=0$.
Thus equality implies $X$ agrees a.e. with a family member.

For a family model, its own index is an available cover at every world,
so $c_i(D_{K_p})\ge p_i$; the budget forces equality at every coordinate.
This proves (a). If $X$ weakly dominates $D_{K_p}$, it is covered everywhere
and $c_i(X)\ge p_i$ for all $i$. Equality in the budget then forces
$X\approx D_{K_p}$. This proves undominatedness.

If $E_i(X)\subseteq E_i(Y)$ a.e. and $Y$ is covered at $i$, every cover of
$Y$ covers $X$ and the correction increases by
$\eta\mu(E_i(Y)\setminus E_i(X))$. Hence
\[
 L_\eta(X,i)\le L_\eta(Y,i)
       -\eta\mu(E_i(Y)\setminus E_i(X)).
\]
If both are uncovered, strict monotonicity follows directly from the error
measure; if only $X$ is covered, it follows from the disjoint loss ranges.

Competitiveness follows from the budget: strict improvement over $D_{K_p}$
at every $i\in\supp p$ would require $\sum_{i\in\supp p}c_i(X)>1$.
The comparison sections are
\[
 \{p:L_\eta(D_{K_p},i)\le L_\eta(X,i)\}
 =\{p:p_i\ge1-L_\eta(X,i)\},
\]
which are closed. Finally set $\beta_i=c_i(X)$ at covered worlds and zero
elsewhere. For a nonfamily model $\sum_i\beta_i<1$, and the index
$q=\beta+(1-\sum_i\beta_i)\one/n$ supplies strict dominance everywhere.
Full coverage follows from $K_r=M$, exactly as in
Theorem~\ref{thm:coverage}.
\end{proof}

This loss does not represent the ordinal preorder from the preceding sections,
which has incomparabilities; nor is it asserted to preserve every strict
comparison of that preorder. It is another construction on the same families.
Its dependence on the full model and optimization over covers makes it different
from the additive CPM error losses. The CPM realisation obstruction therefore
does not extend to all real-valued losses.

\begin{thebibliography}{9}
\bibitem{OR} \emph{An Ordinal Refinement by a Slack Envelope}.
Companion note, 7 October 2026. Source file:
\texttt{Ordinal Refinement by Slack Envelopes.tex}.
\bibitem{LR} \emph{Loss Relations for Sets of Desirable Gambles}.
Companion note. Source file:
\texttt{Loss Relations for Sets of Desirable Gambles.tex}.
\end{thebibliography}
\end{document}
